\documentclass[a4paper,11pt]{article}
\usepackage{fontspec}
\setmainfont{NotoSansCJKsc-Regular.otf}[Path=fonts/,BoldFont=NotoSansCJKsc-Regular.otf,BoldFeatures={FakeBold=2}]
\setsansfont{NotoSansCJKsc-Regular.otf}[Path=fonts/,BoldFont=NotoSansCJKsc-Regular.otf,BoldFeatures={FakeBold=2}]
\XeTeXlinebreaklocale "zh"
\XeTeXlinebreakskip=0pt plus 1pt
\renewcommand{\figurename}{图}
\renewcommand{\tablename}{表}
\renewcommand{\refname}{参考文献}
\usepackage[margin=2.1cm]{geometry}
\usepackage{amsmath,amssymb,booktabs,longtable,array,graphicx,xcolor,hyperref}
\hypersetup{colorlinks=true,linkcolor=black,urlcolor=blue,citecolor=black}
\setlength{\parskip}{0.35em}
\setlength{\emergencystretch}{2em}
\newcommand{\pos}[1]{\left[#1\right]_+}
\title{问题二异构无人机运输调度模型与代码改进说明}
\author{}\date{2026年9月24日}
\begin{document}
\maketitle
\vspace{-2em}
本文核对题目 Word 文档、80箱逐箱清单、16个任务节点、运输机型与电池数据及30米DEM，评估原 \texttt{q2\_solver(8).py}，并给出可复现的改进程序。原程序已实现较完整的物理规则和自适应大邻域搜索，但不能证明全局最优。改进采用真实端点地形遍历、多起点路线搜索、候选路线与时间索引整数规划及独立资源复算。下文将题目规则、建模约定和计算限制分别说明。

\section{文献方法与本题的对应关系}
Cheng等\cite{Cheng2020}研究同构无人机的多架次、时间窗配送，使用二指标模型与分支切割算法处理载荷相关凸非线性能耗，核心包括逻辑割与次梯度割。对本题可借鉴的是逐航段剩余载荷计能、显式检查整架次能量和利用优化界评价求解质量。原文的旋翼功率公式不能替代本题规定的等效航程公式，其模型也不能直接证明本题异构机队及有限共享电池调度的最优性。

Phalapanyakoon与Siripongwutikorn\cite{Phala2023}建立异构、多架次、充电与载荷约束的混合整数模型，目标包含租用、电费和任务超时成本，使用包含机型与路线两部分的染色体、两类初始化、交叉和变异求解。可借鉴异构匹配和路线与充电联动。本题使用固定库存、逐箱时限及独立共享电池，且每次复用前须充满，因此不沿用其机队租金、按剩余电量触发充电或任务整体超时罚金。

将ALNS换成GA并不自动提高质量，也不会产生最优性证明。本次采用混合方法：ALNS产生批次和访问顺序，整数规划在候选池内同时决定保留哪些批次、使用哪种机型以及何时开始。论文名称含有“精确算法”，不代表将其思想移植后的程序也是原问题的精确算法。

\section{原代码核查结果}
\begin{longtable}{p{0.22\textwidth}p{0.72\textwidth}}
\toprule 检查项 & 结论与处理\\\midrule\endhead
非线性能耗与载荷 & 原代码正确按投送后的剩余载荷计算后续航段，返航为空载；保留。\\
医疗与首批硬时限 & 原代码已将医疗期望时间和首批截止时间作为硬要求；两者同时适用时取较早值。\\
共享电池 & 原代码已按独立电池编号和两阶段充电管理，无需将其误判为缺失功能。\\
DEM文件与端点 & 原代码仅读取MAT。本次增加GeoTIFF；按文件的PixelIsPoint标识将tiepoint解释为像元中心，避免额外半像元偏移。原端点吸附算法在120条无向航段中有24条最高高程不同，最大低估14.7147米，最大高估11.8738米。\\
访问点数量 & 原代码将4点上限写入可行性判定，但题目没有该限制。改进允许最多15个不同服务区；大于5点时不再进行阶乘枚举。\\
调度判断 & 单一贪心解码失败不等于原问题无解。增加多个解码规则与整数规划重排，失败时保留已有可行解。\\
修复算子 & 原代码regret相同且代理增量相同后使用负优先系数，破同分方向与注释相反；改为较高优先系数、较紧迫时限优先。\\
可复现性 & 对货箱集合显式排序，避免仅固定随机种子但仍受Python字符串散列顺序影响。记录种子、输入哈希与求解设置。\\
可行性证据 & 从输入货箱和实际路线重新计算载荷、时间、能耗、SOC和充电，检查具体无人机与电池区间冲突，不以搜索器的feasible标记代替验证。\\
输出与路径 & 取消固定Windows路径及必须存在结果模板的要求。自动发现文件，存在同格式多版本时明确报错；输出逐箱、架次、航段、停点、电池CSV及JSON。\\
\bottomrule
\end{longtable}

\section{场景与计算约定}
节点集合为 $V=\{0,1,\ldots,15\}$，其中 $0$ 表示O01。货箱集合 $\mathcal B$ 含80箱。箱 $b$ 的目的地、质量、体积、优先系数和期望时刻分别为 $i_b,w_b,v_b,\omega_b,d_b$。全部物资从O01出发，箱不可拆分且仅交付一次。机型集合 $G=\{A,B,C\}$。
\begin{center}
\begin{tabular}{lrrrrr}\toprule
机型 & 实体无人机 & 共享电池 & 最大载荷/kg & 装载体积/m$^3$ & 完全充电/s\\\midrule
A & 4 & 6 & 25 & 0.060 & 1800\\
B & 2 & 4 & 30 & 0.073 & 2400\\
C & 2 & 4 & 80 & 0.250 & 3000\\\bottomrule
\end{tabular}
\end{center}
全部电池初始满电，返航SOC下限为20\%。共享电池数包含在机电池与备用电池，不能再额外增加8组。不同机型不能混用电池，不同电池允许并行充电，题目未给有限充电桩或工位数量，因此不自行添加数量约束。问题二不使用通信链路和中继数据，也不引入道路通行或村落补点。

本实现约定从准备装载开始，同时占用实体无人机与满电电池；电池返航后立即进入充电。同一停点的货箱统一以该停点全部交接完成时刻记交付，因而是保守交付口径。时间窗只限制送达，不要求返航早于该箱截止。最终完工时间不含最后一批电池的任务后充电。

候选架次采用不重复服务区的基本路线。题目没有给出必须采用该类路线的最优性结论，因此这是当前搜索空间的限制；程序不声称覆盖允许任意重复访问的所有路线。若未来需要允许重复访问，应同时调整停点交付分配与载荷递推，不能只提高点数上限。

\section{航段与架次精确评价}
\subsection{地形净空和航段时间}
设 $\mathcal C_{ij}$ 为由真实端点连接的水平线段经过或接触的DEM像元集合，$z_c$ 为像元高程。采用包含角点接触像元的supercover口径，定义
\begin{equation}
 H_{ij}=\max_{c\in\mathcal C_{ij}}z_c+50,\quad
 h_i=\begin{cases}z_i,&i=0,\\z_i+30,&i\ne0,\end{cases}\quad
 h_{ij}^{+}=H_{ij}-h_i,\quad h_{ij}^{-}=H_{ij}-h_j.
\end{equation}
节点作业高度以附件节点表为准。所有任务航段均验证巡航高度不低于两端作业高度；NoData不能直接略过。程序以球面距离计算小区域水平长度 $d_{ij}$，并在规则经纬度栅格上按实际端点线段遍历。机型 $g$ 的飞行时间为
\begin{equation}
 \tau_{gij}=\frac{h_{ij}^{+}}{v_g^\uparrow}+\frac{d_{ij}}{v_g^c}+\frac{h_{ij}^{-}}{v_g^\downarrow}.
\end{equation}
每次投送后均从服务区30米作业高度重新爬升。不能用整条多点路线的统一巡航高度替换题目规定的逐航段高度。

\subsection{非线性能耗和后缀载荷}
机型 $g$ 的最大载荷、空满载标准航程、可用电量、含电池空载质量和爬升效率分别记为 $Q_g$、$L_g^0$、$L_g^F$、$E_g^{use}$、$m_g^0$ 和 $\eta_g$。沿用题目等效航程模型，并按标准航程定义水平耗能：
\begin{align}
 L_g(q)&=L_g^0-(L_g^0-L_g^F)\left(\frac q{Q_g}\right)^{3/2},\\
 E_{gij}(q)&=E_g^{use}\frac{d_{ij}}{L_g(q)}
       +\frac{(m_g^0+q)\,9.81\,h_{ij}^{+}}{3.6\times10^6\eta_g}.
\end{align}
能量单位为kWh，时间为秒。含电池空载质量已经包含电池，不重复计入。下降附加能耗按附件取零；题目没有提供运输悬停功率，本次不另加运输悬停能耗。

架次 $p=(0,i_1,\ldots,i_m,0)$ 在停点 $i_l$ 交付箱集 $B_{pl}$。进入该停点前的载荷为
\begin{equation}
 q_{pl}=\sum_{k=l}^{m}\sum_{b\in B_{pk}}w_b,\qquad q_{p,m+1}=0.
\end{equation}
设 $B_p=\bigcup_l B_{pl}$。可行架次必须满足
\begin{equation}
 \sum_{b\in B_p}w_b\le Q_g,\qquad
 \sum_{b\in B_p}v_b\le V_g,\qquad
 E_p=\sum_{l=1}^{m+1}E_{g,i_{l-1},i_l}(q_{pl})\le(1-\rho_g)E_g^{use},
\end{equation}
其中 $i_0=i_{m+1}=0$。质量与体积均随投送单调减少，故在初始总载荷处检查容量即可。不同访问顺序具有不同后缀载荷与能耗，不能只按最短距离排序。

\subsection{逐箱交付时刻与硬时限}
记架次准备时间 $P_p=t_g^{prep}+|B_p|t_g^{load}$，第 $l$ 个停点交接时间 $u_{pl}=t_g^{base}+|B_{pl}|t_g^{box}$。从准备开始到停点交接结束的偏移为
\begin{equation}
 \delta_{pl}=P_p+\sum_{k=1}^{l}\big(\tau_{g,i_{k-1},i_k}+u_{pk}\big),\qquad
 D_p=\delta_{pm}+\tau_{g,i_m,0}.
\end{equation}
若架次从 $s_p$ 开始，则 $b\in B_{pl}$ 的送达时刻是 $T_b=s_p+\delta_{pl}$，返航时刻为 $F_p=s_p+D_p$。定义
\begin{equation}
 \bar d_b=\min\left(\{d_b:b\text{为医疗箱}\}\cup
 \{d_b^{first}:b\text{为首批箱}\}\right),
\end{equation}
空集合的最小值记为 $+\infty$。硬约束为 $T_b\le\bar d_b$。其他箱的期望时刻进入及时性指标，不强制全部作为硬时间窗。架次最晚开始时刻为
\begin{equation}
 \ell_p=\min_{b\in B_p}(\bar d_b-\delta_{bp}).
\end{equation}

\subsection{电池周转}
返航SOC与充电时长分别为 $\sigma_p=1-E_p/E_g^{use}$ 及
\begin{equation}
 c_p=T_g^{full}\begin{cases}
 0.65(0.90-\sigma_p)/0.90+0.35,&0\le\sigma_p<0.90,\\
 0.35(1-\sigma_p)/0.10,&0.90\le\sigma_p\le1.
 \end{cases}
\end{equation}
同一无人机的任务区间 $[s_p,F_p)$ 不得重叠；同一电池的任务加充电区间 $[s_p,F_p+c_p)$ 不得重叠。电池可在相同机型的不同无人机间流转，无需与某一无人机终身绑定。

\section{多目标与联合调度模型}
配送及时性、完工时间、能耗与架次数定义为
\begin{equation}
 L=\sum_{b\in\mathcal B}\omega_b\frac{\pos{T_b-d_b}}{d_b},\quad
 C_{\max}=\max_p F_p,\quad E=\sum_p E_p,\quad N=\text{执行架次数}.
\end{equation}
医疗与首批要求始终是硬约束，不能通过降低罚系数容忍违反。$L$ 是无量纲加权相对迟到，不是迟到秒数。本题 $d_b>0$。归一化尺度固定为同一初始可行解：$L_0=\max(1,L^{init})$，其余尺度为 $C_0,E_0,N_0$，所有随机种子和偏好方案共用。优化目标为
\begin{equation}\label{eq:objective}
 \min J=\alpha_L\frac L{L_0}+\alpha_C\frac{C_{\max}}{C_0}
          +\alpha_E\frac E{E_0}+\alpha_N\frac N{N_0}.
\end{equation}
权重表示建模偏好，而非题目给定常数。程序计算四组偏好，并对结果集做非支配性检查：
\begin{center}\begin{tabular}{lrrrr}\toprule
偏好 &$\alpha_L$&$\alpha_C$&$\alpha_E$&$\alpha_N$\\\midrule
均衡 &0.45&0.25&0.20&0.10\\
完工时间 &0.45&0.40&0.10&0.05\\
能耗 &0.45&0.10&0.40&0.05\\
架次数 &0.45&0.10&0.10&0.35\\\bottomrule
\end{tabular}\end{center}
不同权重可能选中同一个方案，这表示当前候选解集中不存在足以改变选择的权衡，不能伪造四套不同路线。权重和并不能找到全部非凸Pareto前沿。需要更明确的业务偏好时，可改用 $L\le\varepsilon_L$、$C_{\max}\le\varepsilon_C$ 等约束再最小化能耗或架次。

\subsection{候选路线与时间索引整数规划}
每个候选 $p\in\mathcal P$ 已包含具体箱集、访问顺序和机型 $g(p)$，因此非线性能耗已按原式算成常数。变量 $x_{pk}\in\{0,1\}$ 表示候选在时刻 $k\Delta$ 开始。仅生成满足 $k\Delta\le\ell_p$ 且 $k\Delta+D_p\le H$ 的列。记 $a_{bp}=1$ 表示候选包含箱 $b$，$n_g$、$B_g$ 分别为机型无人机数和共享电池数。

逐箱集合划分约束为
\begin{equation}
 \sum_{p\in\mathcal P}\sum_k a_{bp}x_{pk}=1,\qquad b\in\mathcal B.
\end{equation}
资源占用向上取整，$d_p^\Delta=\lceil D_p/\Delta\rceil$，$b_p^\Delta=\lceil(D_p+c_p)/\Delta\rceil$。每个时间槽 $t$ 满足
\begin{align}
 \sum_{p:g(p)=g}\sum_{k:k\le t<k+d_p^\Delta}x_{pk}&\le n_g,\\
 \sum_{p:g(p)=g}\sum_{k:k\le t<k+b_p^\Delta}x_{pk}&\le B_g.
\end{align}
由于开始时刻均落在格点上，这种占用检查不会漏掉格内真实重叠。对每个箱加入
\begin{equation}
 C_{\max}\ge\sum_{p,k}a_{bp}(k\Delta+D_p)x_{pk}.
\end{equation}
每个箱恰好属于一个选中架次，因此这些约束的共同最大值正是最晚返航时刻。还加入机型总工作量下界 $n_gC_{\max}\ge\sum_{p:g(p)=g,k}D_px_{pk}$ 以加强松弛。及时性项可预计算为
\begin{equation}
 L_{pk}=\sum_{b\in B_p}\omega_b\frac{\pos{k\Delta+\delta_{bp}-d_b}}{d_b},\qquad
 L=\sum_{p,k}L_{pk}x_{pk}.
\end{equation}
于是式\eqref{eq:objective}是线性目标，所有物理能耗仍为非线性原式的精确评价结果，并没有用线性近似代替能耗。

\subsection{从资源数量到具体编号}
同一机型内的无人机和电池参数一致，初始均可用。给定选中的时间区间后，按开始时刻排序，将每个区间分配给最早释放的实体。区间图的染色数等于最大同时占用数，故上述两组容量约束足以恢复有限库存内的具体无人机和电池编号。无人机占用区间与电池占用区间分别染色，再按架次配对即可；这利用了电池可在同机型实体间共享的条件。

\section{求解流程与最优性边界}
首先按医疗和首批任务构造紧急种子，普通物资按容量和物理能量形成批次；每轮ALNS执行随机、地理相关、迟到相关或整条路线移除，再进行regret-2修复。多个调度规则分别尝试机型与时序安排，以真实完整调度指标择优。保存搜索中出现的物理可行路线，并进行两个整数规划阶段：固定箱集与访问顺序的机型时序优化，以及允许候选批次重组的联合优化。

整数规划完成后，保持每个具体资源上的架次先后关系，以连续时间向前压紧空隙。该操作不会推迟任何交付、改变能耗或增加架次，因此不会恶化目标。求解器时限结束无可行整数解时，保留已验证的启发式方案。

本次时间网格为60秒，计算时域由初始完工时间的1.4倍确定，路线池为有限搜索所得。求解器报告的下界和gap只属于其记录的候选池、时间网格及计算时域。即使gap为零，也只能称该受限模型求解完成；它不是原连续时间问题的全局下界或全局最优证明。压紧后的解可能离开原开始时刻网格，不能再直接套用压紧前的gap。要证明原问题全局最优，需要完整可行路线空间与相应的精确定价、分支定价/切割，或可证完整的连续时间混合整数模型。

\section{附件算例结果与资源验证}
所有数值均来自本次附件实跑。基线保留原ALNS与原解码器，统一更正地形模块，运行800次迭代、种子20260924、固定Python散列种子0。改进运行四组偏好，每组两个种子、每个500次迭代，并以基线路线热启动，随后求解有限时长的整数规划。因此下表证明本算例取得了更好的可行方案，不是相同计算预算下的算法排名。
\begin{center}
\begin{tabular}{lrrrrr}\toprule
方案 & 架次 & 完工/min & 能耗/kWh & 加权迟到 & 按期箱数\\\midrule
原ALNS统一地形 & 20 & 159.72 & 64.5185 & 6.6369 & 77/80 \\
改进均衡方案 & 20 & 137.48 & 61.8969 & 0.0000 & 80/80 \\
减少架次方案 & 17 & 155.22 & 60.6636 & 7.2047 & 75/80 \\
\bottomrule\end{tabular}\end{center}
均衡方案比基线缩短完工时间22.24分钟（13.92\%），减少能耗2.6216kWh（4.06\%），架次数相同，全部80箱在期望时间内交付。均衡、完工时间和能耗三组权重最终选择了同一方案。减少架次方案使用17架次，5箱普通物资迟于期望时刻，但全部医疗和首批箱仍满足硬时限。

均衡方案的31个不同硬时限箱（16箱医疗、30箱首批，两组有重叠）全部满足要求。最紧硬时限余量为206.804秒，最低返航SOC为20.971\%，高于20\%下限。每箱恰好一次；具体无人机无任务重叠；具体电池的任务和充电区间无重叠。独立审计使用微秒级端点归一化处理浮点累计误差，不把同一资源前后任务的共同端点计为重叠。
\begin{center}\begin{tabular}{lrrrr}\toprule
机型 & 同时执行无人机峰值 & 库存 & 电池任务加充电峰值 & 电池库存\\\midrule
A & 4 & 4 & 4 & 6 \\
B & 2 & 2 & 4 & 4 \\
C & 2 & 2 & 4 & 4 \\
\bottomrule\end{tabular}\end{center}
独立地形验证将全部120条无向航段与另一套线段矩形相交计算交叉核对，同时检查穿角点、沿网格边界与90\%充电分段边界。修改已导出结果中的能耗、无人机ID或电池ID，校验器能够识别错误。

本次实际增益来自多起点路线搜索及固定路线的机型时序整数规划，后者将均衡ALNS解的完工时间从8909.39秒进一步降至8248.92秒。四次较大的候选池联合模型均在当前时限内未取得可用整数解，因而没有用它们替换已验证方案。程序完整实现该阶段，但不能将“已建立模型”描述为“已取得更优联合解”。固定路线子问题的求解状态如下，均为压紧前的离散模型：
\begin{center}\begin{tabular}{lrrr}\toprule
偏好 & 开始时间列数 & 相对gap & 实际求解/s\\\midrule
均衡 & 3745 & 0.0000\% & 4.02 \\
完工时间 & 4042 & 0.0299\% & 5.80 \\
能耗 & 3130 & 0.3518\% & 1.18 \\
架次 & 2729 & 0.0000\% & 1.57 \\
\bottomrule\end{tabular}\end{center}
均衡固定路线模型的gap为零，但其箱集、顺序和开始网格均受到限制，因此不能称原问题全局最优。最晚返航的统计口径严格为所有运输无人机完成最后一个架次并返回O01的最大时刻。


应急及时性与节能、减少架次通常存在权衡：合并多箱和多服务区可能节省固定准备及往返开销，却增加初始载荷和后续货箱等待；提前派出轻载批次有利于紧迫物资，但可能增加架次。大机型载荷大，不能据此认定单位任务能耗最低。更多电池只在充电周转约束紧张时有帮助，并不能突破实体无人机数。

\begin{figure}[htbp]\centering
\includegraphics[width=.98\linewidth]{results/drone_schedule.pdf}
\caption{均衡方案具体无人机调度，浅色段表示准备装载。横轴以分钟显示，计算仍采用秒。}
\end{figure}
\begin{figure}[htbp]\centering
\includegraphics[width=.98\linewidth]{results/battery_schedule.pdf}
\caption{均衡方案共享电池使用与充电。斜线区间为充电，最后一次充电不计入任务完工时间。}
\end{figure}

\section{代码与结果使用}
程序只读取节点、货箱、运输无人机和DEM四类输入。压缩包已包含本次输入和已计算的各偏好方案。命令行示例：
\begin{verbatim}
python q2_solver_improved.py --root data --output new_results 
  --iters 500 --seeds 20260924 20260925 --milp-seconds 90
python verify_q2.py --root data
  --solution results/balanced/solution.json
\end{verbatim}
实际输入命令时将续行合并为一行，Windows用户也可使用随附批处理文件。要复现本次有基线热启动的运行，添加
\texttt{--warm-routes baseline\_routes.json}。若不提供，该文件不是必需输入。

时间上限可能导致不同CPU上整数规划取得不同的当前最好解，已保存的逐架次计划及独立校验不依赖重跑搜索。

CSV采用UTF-8带BOM，可直接在Excel打开。各偏好文件夹分别保存架次、逐箱、航段、交接与电池表，另有机器可读JSON和审计报告。根目录保存输入哈希、种子、指标尺度、各子问题状态和局部gap。当前Excel汇总为本次已计算结果，修改权重需要重新运行优化程序；表内重新计算不能替代路线优化。模型LaTeX可直接用XeLaTeX编译，论文正文可选取所需段落并调整章节层级。

\begin{thebibliography}{9}
\bibitem{Cheng2020} Cheng C, Adulyasak Y, Rousseau L M. Drone routing with energy function: Formulation and exact algorithm[J]. Transportation Research Part B: Methodological, 2020, 139: 364--387. DOI: \href{https://doi.org/10.1016/j.trb.2020.06.011}{10.1016/j.trb.2020.06.011}.
\bibitem{Phala2023} Phalapanyakoon K, Siripongwutikorn P. Route planning of heterogeneous unmanned aerial vehicles under recharging and mission time with carrying payload constraints[J]. Journal of Industrial Engineering and Management, 2023, 16(2): 215--235. DOI: \href{https://doi.org/10.3926/jiem.4381}{10.3926/jiem.4381}.
\end{thebibliography}
\end{document}
